\section*{Response to the Editor\bigskip{}
}

Dear Editor,\bigskip{}

Thank you for the opportunity to revise and resubmit our paper. We appreciate the thoughtful comments from the referees and have addressed their concerns in this revision.\medskip{}

% Add your response to the Editor here

%%%% DS: I made changes to the next two paragraphs. Original version below.
We now frame the paper as a frictionless benchmark model in which beliefs serve as a direct source of business cycle fluctuations. 
In this benchmark, there are no financial frictions, no short-selling constraints, prices are fully flexible, and markets are complete. 
Beliefs and output are only linked through a timing assumption: firms must choose labor input before observing their productivity. 
We see value in developing such a frictionless benchmark, as it allows us to isolate the direct effects of beliefs on business cycles. 
These effects are likely to be present-- and potentially be amplified---in environments that feature additional frictions. \medskip{}

% We also invested effort to produce new theoretical insights and provide greater empirical support. 
While your second route suggested we focus on the qualitative properties, we also pursued a quantitative evaluation of the model. 
% Our goal was to show the extent the model can produce quantitatively sensible predictions.  
We purposely avoid introducing many bells and whistles, as we wanted to keep the model parsimonious, introducing only minor changes to facilitate the calibration.
Despite its simplicity, the model produces reasonable predictions for both asset prices and labor market fluctuations. 
Section~\ref{sec:quantitative} now includes a more comprehensive quantitative assessment, which substantially improves upon the earlier version we had presented. \medskip{}
